\begin{proposition}[Normal crossing torus fibres and their normalizations]\hypertarget{tag-02Y0}{}\label{night:P4.8}\label{stmt:P4.8}
\StatementRecord{night:P4.8}{proved}{02Y0}{\SeriesRef{night:P4.12}, \SeriesRef{night:P4.6}, \SeriesRef{night:P4.7}}

For a fibration in Proposition~\ref{night:P4.7} with general fibre
a complex two-torus, an irreducible singular support with
rational normalization has Euler number two; every other irreducible
support has Euler number zero. The supports of Euler number zero have the following forms: a smooth support has a finite unramified
torus cover; a singular irreducible support has smooth minimal
elliptic-ruled or Hopf normalization; a reducible support has smooth
minimal Hopf or non-Kähler elliptic normalization components, with
smooth elliptic double curves and finite unramified branch maps.
Thus $c_3[X]=0$ precisely when every fibre belongs to these zero cases.
\end{proposition}
\begin{proof}

\begin{enumerate}[label={$\langle1\rangle$\arabic*.},leftmargin=*]
\item We apply \cite[Chapter~VI, Theorem~1.1 and Table~10, pp.~243--244]{BHPV2004E} to the smooth compact
normalization components, using the identity $e(S_i)=K_{S_i}^2$ proved
in Proposition~\ref{night:P4.6}. The relevant non-Kähler classification is stated in
\cite[arXiv version~3, Theorem~4.2, pp.~17--18; Theorem~5.1 and proof, pp.~27--29]{OVV2020}: a minimal non-Kähler
surface is elliptic or of class VII, and the reduced fibres of a
minimal non-Kähler elliptic surface are smooth elliptic curves.

\item For a component with nonempty conductor, a Kähler form $\omega$ would
give $K_{S_i}\cdot\omega=-D_i\cdot\omega<0$. No positive power of
the canonical bundle can then have a nonzero section. The classification
makes a Kähler component rational or ruled. If it is obtained by $k$
blowups of a ruled surface over a genus-$g$ curve, then
\[
 e=4(1-g)+k,\qquad K^2=8(1-g)-k.
\]
Equality gives $k=2(1-g)$, so $g\le1$ and the genus-one case is minimal.
A rational component has $\chi(\OO)=1$, hence $e=6$ here, including
the case with minimal model $\Pp^2$.

\item A minimal non-Kähler elliptic surface has $e=K^2=0$, by its smooth
elliptic reduced fibres. A minimal class-VII surface has $b_1=1$ and
$p_g=0$, so $\chi(\OO)=0$ and $K^2=-e$, with $e=b_2\ge0$.
Blowups preserve $\chi(\OO)$ and change $(e,K^2)$ by $(1,-1)$.
Consequently $e=K^2$ forces minimality and $e=0$ in all these
non-Kähler cases. The class-VII case has $b_2=0$. The finite-cover classification
\cite[arXiv version~3, Theorem~1, p.~1]{BBK2019} gives a finite
unramified map $p:\widehat S_i\to S_i$ with $\widehat S_i$ Hopf
or Inoue. The inverse image $p^{-1}(D_i)$ is a nonempty compact
curve: surjectivity makes it nonempty, and an unramified map preserves
its local codimension. Inoue surfaces have no compact curves
\cite[Theorem~3.1, p.~303]{KhaledTeleman2026Journal}; the type-III
family is covered by the explicit unramified type-II double cover
of \cite[proof of Remark~3.5, p.~305]{KhaledTeleman2026Journal}.
Thus $\widehat S_i$ is Hopf. Its universal covering
$\C^2\setminus\{0\}\to\widehat S_i$, composed with $p$, is a
universal covering of $S_i$, which is therefore a Hopf surface.

\item For an irreducible singular support there is one normalization
component. Proposition~\ref{night:P4.6} gives $T=2\chi(\OO_{S_1})$, which is two in the
rational case and zero in the other cases. Proposition~\ref{night:P4.12} proves that the
dualizing sheaf of an irreducible support is torsion. Adjunction on
the normalization therefore makes $K_{S_1}+D_1$ torsion. Since
$D_1\ne0$, every positive plurigenus of $S_1$ vanishes: multiplying
a pluricanonical section by an appropriate positive power of the
conductor section would give a section of a torsion line bundle with
a zero, whose sufficiently high power is a nonzero constant. Thus a
non-Kähler normalization in this irreducible case is of class VII,
and hence is the Hopf case already obtained.

\item For a reducible connected fibre, each normalization component meets
another global irreducible component. Pulling back the divisor of
that component gives a nonzero effective divisor on $S_i$ with zero
real first Chern class. A Kähler form would pair positively with it,
contradicting this zero class. Hence every component is non-Kähler,
has $\chi(\OO)=0$, and $T=0$. The conductor components are disjoint
and smooth, because all their nodes would come from triple points.
For any conductor component $C$, adjunction and $K+D\equiv0$ give
$2g(C)-2=(K+C)\cdot C=0$, since the other conductor components
are disjoint. Thus the double curves are smooth elliptic, and their
branch maps are the unramified maps already described in Proposition~\ref{night:P4.6}.

\item Finally let the support $G$ of a multiple fibre $mG$ be smooth. The
normalized cyclic base change $t=u^m$ is locally a disjoint union of
smooth charts: the original equation is $t=z^m$ times a holomorphic
unit, whose root can be extracted locally. These charts show that
the central fibre of the base-changed family maps finitely and
unramifiedly to $G$, and that its map to the $u$-disc is a proper submersion. Its nearby
fibres are the original tori. Catanese's theorem
\cite[definition of deformation equivalence, p.~261; Theorem~2.1 and proof, pp.~264--268]{Catanese2004} states that every
deformation of a complex torus is a complex torus; its definition
uses precisely a proper holomorphic submersion over a connected
disc. The normalized cyclic base change constructed above is such
a proper holomorphic submersion, so its central fibre is a torus;
Ehresmann's theorem also shows that it is connected. This gives
the finite unramified torus cover of $G$.

\end{enumerate}
\end{proof}